%%%%%%%%%%%%%%%%%%%%%%%%%%%%% Define Article %%%%%%%%%%%%%%%%%%%%%%%%%%%%%%%%%%
\documentclass{article}
%%%%%%%%%%%%%%%%%%%%%%%%%%%%%%%%%%%%%%%%%%%%%%%%%%%%%%%%%%%%%%%%%%%%%%%%%%%%%%%

%%%%%%%%%%%%%%%%%%%%%%%%%%%%% Using Packages %%%%%%%%%%%%%%%%%%%%%%%%%%%%%%%%%%
\usepackage{listings, xcolor}
\usepackage{parskip}
\usepackage{hyperref}
\usepackage{amsmath}
\usepackage{graphicx}
\usepackage[a4paper, margin = 1.5cm]{geometry}
\usepackage{float}
\usepackage{amssymb}

%%%%%%%%%%%%%%%%%%%%%%%%%%%%%%% Title & Author %%%%%%%%%%%%%%%%%%%%%%%%%%%%%%%%
\title{Esercizi BDD}
\author{Trentin Tommaso}
%%%%%%%%%%%%%%%%%%%%%%%%%%%%%%%%%%%%%%%%%%%%%%%%%%%%%%%%%%%%%%%%%%%%%%%%%%%%%%%

\begin{document}
  \section{Esercitazione 1}
    \subsection{Es 1}
      Si supponga di realizzare una rete a commutazione di pacchetti usando una tecnologia di collegamento punto-punto che garantisce consegna affidabile a livello DL, a meno di guasti completi del collegamento:
      \begin{enumerate}
        \item Cosa può fare un nodo se si accorge che un collegamento è guasto?
        \item Necessario implementare altri meccanismi di gestione errori ai livelli superiori (es. trasporto)?
      \end{enumerate}
      \begin{enumerate}
        \item R - In caso di guasto si cambia "percorso" per arrivare a destinazione
        \item R - Necessario a causa del possibile arrivo out-of-order dei pachhetti
      \end{enumerate}
    \subsection{Es 2}
      La lingua italiana conta 45 fonemi, in un normale eloquio vengono pronunciati 15 fonemi al secondo:
      \begin{enumerate}
        \item A quanto equivale il bitrate di tale trasmissione?
        \item Se il canale usato per la trasmissione ha SNR di 10dB, quale bandwidth è necessaria?
      \end{enumerate}
      \begin{enumerate}
        \item R - $\log_2(45)*15=\sim82.38$
        \item R - $C=B\log_2(1+\frac{S}{N}) \rightarrow 82.38=B\log_2(1+10) \rightarrow B=82.38/\log_2(11)=23.81 [Hz]$  
      \end{enumerate}
    \subsection{Es 3}
      Nei sistemi radio DAB e DAB+ i dati sono codificati mediante simboli (OFDM): ogni simbolo porta 3072 bit e dura 1.246 ms. Ogni frame contiene 76 simboli e i frame sono separati da una pausa di 1.304 ms
      \begin{enumerate}
        \item A quanto equivale il bitrate grezzo di tale trasmissione?
        \item Se il canale usato ha una bandwidth di 1536 kHz, quanto dev'essere il rapporto S/N minimo per una trasmissione senza errori?
      \end{enumerate}
      \begin{enumerate}
        \item R - $76*1.246+1.304=96 [ms] \rightarrow \frac{76*3072}{96}=2.432 [Mbps]$
        \item R - $C=B\log_2(1+\frac{S}{N}) \rightarrow \log_2(1+\frac{S}{N})=\frac{2.432}{1536} \rightarrow 1+\frac{S}{N}=2^(\frac{2.432}{1536})=3 [dB]$
      \end{enumerate}
    \subsection{Es 4}
      Un canale di trasmissione ha bandwidth di 10 MHz e un rapporto S/N in potenza di -20 dB:
      \begin{enumerate}
        \item Determinare massimo bitrate "grezzo" di tale canale
        \item Se si usa codifica con un alfabeto di 8 simboli, qual'è il baudrate corrispondente?
      \end{enumerate}
      \begin{enumerate}
        \item R - $SNR=10\log(\frac{S}{N}) \rightarrow \frac{S}{N}=10^(\frac{SNR}{10}) \rightarrow C=10MHz\log_2(1+10^(-2))=0.143 [Mbps] = 143 [Kbps]$
        \item R - $\frac{143}{3}=47 [Kbaud]$
      \end{enumerate}
    \subsection{Es 5}
      Controllo missione NASA sta comunicando con rover su Marte, che dista 60 milioni di Km dalla Terra, su un canale con capacità 250 Kbps
      \begin{enumerate}
        \item Ricordando che $c=3*10^8 \frac{m}{s}$, quant'è il RTT?
        \item Periodicamente il rover scatta una foto che pesa 5MB e la trasmette su tale canale, dopo quanto tempo dal momento dello scatto quella foto è disponibile presso il controllo missione?
      \end{enumerate}
      \begin{enumerate}
        \item R - $\frac{60*10^9}{3*10^8}=20*10=200s \rightarrow RTT=200*2=400 [s]$
        \item R - $5MB=40Mb \rightarrow \frac{40000 Kb}{250 \frac{Kb}{s}}=160 [s] \rightarrow 200+160=360 [s]$
      \end{enumerate}
    \subsection{Es 6}
      Il sistema DVB-T2 prevede diversi profili a livello fisico. Uno di questi definisce una modulazione 64-QAM (alfabeto di 64 simboli) con una FEC 3/5, su un canale largo 8 MHz; si codifica un simbolo per ciclo. Sapendo che circa l'8\% della banda viene consumato da traffico di servizio, quant'è il bitrate per ogni canale? \\
      R - $8MHz*1 \frac{sym}{cyc} = 8 [Mbaud] \rightarrow 8 \frac{Mbaud}{s} * 6 \frac{bit}{sym} = 48 [Mbps] \rightarrow 48Mbps*\frac{3}{5}= 28.8 [Mbps] \rightarrow 28.8*\frac{92}{100}=26.5 [Mbps]$
    \subsection{Es 7}
      Su una certa linea di trasmissione i dati vengono inviati in frame composti da una intestazione di 10 byte comprensiva di CRC e un payload di 1000 byte. Se un frame è errato viene scartato e ritrasmesso. La probabilità di errore per ogni bit è $p=10^(-5)$. Se questi frame vengono usati per trasferire un file da 10 KB qual'è la probabilità che si debba effettuare almeno una trasmissione? \\
      R - $1000B+10B=1010B=8080 [b] \rightarrow$ probabilità trasmissione corretta= $(1-p)^8080=0.922 \rightarrow$ probabilità ritrasmissione= $1-ptc=1-0.922^10=~5\%$ 
    \subsection{Es 8}
      Su una linea con probabilità di errore per bit di 10^(-4) si trasmettono due frame, ognuno di 500 byte. Qual'è la probabilità che almeno uno arrivi a destinazione senza errori? \\
      R - Probabilità un frame con errori= $1-((1-10^(-4))^(500*8))=0.3297$, probabilià almeno uno senza errori: $1-(0.3297^2)=0.8913$
\end{document}
